\section{Temporal Contracts and Evidence Requirements}\label{sec:setup}

At observation $o_t$ with instruction $l$, a VLA generates $H$ valid actions
\begin{equation}
  A_t=(a_t,\ldots,a_{t+H-1}),
\end{equation}
and executes $E\leq H$ actions before the next query.  Let $\tau(t)$ denote the
frame represented by a future condition $z_t^*=F(o_{\tau(t)})$, and let
$\hat z_t$ be its prediction.  We summarize the target horizon relative to the
model and execution endpoints as
\begin{equation}
 g_t^H=\frac{\tau(t)-t}{H-1},\qquad
 g_t^E=\frac{\tau(t)-t}{E-1}.
\end{equation}
A model-endpoint target has $g_t^H=1$; an execution-endpoint target has
$g_t^E=1$; and a multi-query milestone has $g_t^E>1$.  These definitions
describe timing, not usefulness.  A multi-query target can be effective when a
planner, time-to-go input, or learned hierarchy translates it into a local
action constraint.

The released RoboTwin LaWAM contract has $H=E=36$ and samples its second visual
frame at offset 35.  Its target therefore coincides with both endpoints under
the audited current two-frame loader.  The historical local LMWM contract has
$H=50$ and $E=36$, so its fixed target at offset 49 lies beyond the executed
prefix.  Recurrence milestones vary more strongly: across 420,238 frozen pairs,
their mean horizon is 117.45 frames, only 19.35\% fall within the 50-action
model window, and 12.50\% fall within the executed 36-action prefix.  The
released training stream is unavailable, so endpoint alignment is inferred
from the released config and audited loader rather than reconstructed examples.

We require four levels of evidence:
\begin{enumerate}
  \item \textbf{Predictability:} $\hat z_t$ improves a held-out representation
  metric over persistence or shuffled actions.
  \item \textbf{Matched utility:} independently trained predictive and control
  policies differ in closed-loop success across training seeds.
  \item \textbf{Content use:} at a fixed checkpoint, correct $\hat z_t$
  outperforms paired shuffled content; null and persistence diagnose route and
  future-state dependence separately.
  \item \textbf{Source attribution:} matched component arms isolate
  pretraining, auxiliary shaping, inference conditioning, and their interaction.
\end{enumerate}
No level implies the next.  In particular, representation quality cannot
establish utility, a system score cannot identify a component, and simulator
seeds cannot substitute for training-seed replication.
